%% energese-shapes.tex -- Odum Energy Systems Language symbols as PGF shapes.
%%
%% Every shape sizes itself to its label. \energese@measure derives a text
%% half-width and half-height from \pgfnodeparttextbox plus the inner
%% separation, clamped against `minimum width' and `minimum height'; each shape
%% then grows its outline around that rectangle by a shape-specific margin, so
%% the label always sits inside the symbol. Fixed-size symbols clip anything
%% longer than a word.
%%
%% Geometry follows Odum's canonical symbol set:
%%   source       circle
%%   storage      circle with a flattened, peaked top (a tank)
%%   producer     rectangle closed by a semicircle on the right
%%   consumer     hexagon, pointed top and bottom, with vertical sides
%%   transaction  diamond (Odum's "exchange")
%%   interaction  pointed right, notched left
%%   gain         right-pointing triangle (constant-gain amplifier)
%%   loop_limited rectangle whose right side bows outward
%%   switch       square with all four sides bowed inward
%%   ground       fixed-size ground glyph; carries no label
%%
%% Anchor conventions used by the layout engine: energy enters at `energy_in'
%% on the left and leaves at `energy_out' on the right; money moves through
%% `money_in'/`money_out' in the upper half; dissipation leaves through
%% `heat_sink' at the bottom.
%%
%% Proportions are measured from reference-images/complete-odum-symbols.png by
%% flood-filling each symbol's interior, which excludes the attached arrows.
%% tests/fidelity/invariants.json records the targets and tolerances, and
%% `make conformance' checks the rendered symbols against them. Do not adjust
%% these numbers by eye.

% Own scratch registers rather than \pgf@xa and friends, so that an
% intervening \pgfmathsetlength cannot disturb a partial measurement.
\newdimen\energese@dimx
\newdimen\energese@dimy
\newdimen\energese@tmp

% Measure the node's text box. Leaves the half-width in \energese@dimx and the
% half-height in \energese@dimy, each at least half the requested minimum size.
\def\energese@measure{%
    \energese@dimx=.5\wd\pgfnodeparttextbox%
    \pgfmathsetlength\energese@tmp{\pgfkeysvalueof{/pgf/inner xsep}}%
    \advance\energese@dimx by\energese@tmp%
    \energese@dimy=.5\ht\pgfnodeparttextbox%
    \advance\energese@dimy by.5\dp\pgfnodeparttextbox%
    \pgfmathsetlength\energese@tmp{\pgfkeysvalueof{/pgf/inner ysep}}%
    \advance\energese@dimy by\energese@tmp%
    \pgfmathsetlength\energese@tmp{\pgfkeysvalueof{/pgf/minimum width}}%
    \energese@tmp=.5\energese@tmp%
    \ifdim\energese@dimx<\energese@tmp%
        \energese@dimx=\energese@tmp%
    \fi%
    \pgfmathsetlength\energese@tmp{\pgfkeysvalueof{/pgf/minimum height}}%
    \energese@tmp=.5\energese@tmp%
    \ifdim\energese@dimy<\energese@tmp%
        \energese@dimy=\energese@tmp%
    \fi%
}

% Fit an outline to a label without distorting it.
%
% A symbol must hold its shape at every size: widening a hexagon to swallow a
% long label makes it read as a different symbol. So each shape declares its
% canonical aspect R (outline width over height, measured from Odum's sheet)
% and the fractions of that outline its text may occupy, fx and fy. The outline
% is then the smallest size at aspect R which contains the label, and the label
% -- not the shape -- is what gives way when it will not fit.
%
% \energesemagnitude scales the symbol only. Applying TikZ's `scale' to the node
% would scale its label too, so two symbols of different magnitude carrying the
% same words rendered them at different sizes.
%
% Leaves the outline half-width in \energese@outw and half-height in
% \energese@outh.
\newdimen\energese@outw
\newdimen\energese@outh
\def\energesemagnitude{1}

\def\energese@fit#1#2#3{%
    \energese@measure%
    \pgfmathsetlength\energese@tmp{%
        \energesemagnitude * max(\the\energese@dimx/(#2*#1), \the\energese@dimy/#3)}%
    \energese@outh=\energese@tmp%
    \pgfmathsetlength\energese@tmp{#1*\the\energese@outh}%
    \energese@outw=\energese@tmp%
}

% Connection ring -- the symbol's ports.
%
% Every symbol carries a ring of connection points spaced a fixed distance
% apart *along its outline* -- \energeseanchorpitch, 3mm by default. Spacing
% rather than a count is the point: a coupled pair attaching at `ring 3' and
% `ring 4' is separated by the same distance whichever symbol it meets, so
% couplings read alike across a diagram. A large symbol therefore carries more
% ports than a small one. (The pitch is nominal to within one port in
% \energese@ringN below, which explains why.)
%
% These are the ports: a pathway attaches to one of them and to nothing else,
% and a port holds one pathway (see \energeseport in energese.sty). The named
% anchors (energy_in, money_out, ...) name a *direction* -- where a kind of
% flow belongs, as Odum draws it -- and the port used is the ring point nearest
% that direction, which for the anchors on the axes is the same point. A
% connection point free to sit between ports would put a pathway somewhere the
% symbol does not offer one.
%
% The anchors are named by index, `ring 0' at 0 degrees and counting
% anticlockwise, because the angle of a given index depends on the symbol's
% size. Indices run to 47; past the ring's last port they wrap, harmlessly
% repeating earlier positions.
%
% The angular step comes from the pitch over the symbol's effective radius,
% which is exact for a circle and an approximation elsewhere; each shape's
% radius-at-angle macro then puts the point on its actual outline. Being *on
% the outline* is the part that must be exact -- a port off the glyph is a
% pathway attached to nothing -- and every radius macro below is exact, or
% converged to better than 0.05%, for that reason.
\newlength{\energeseanchorpitch}
\setlength{\energeseanchorpitch}{3mm}

% How many ports fit around this outline: the pitch sets the angular step, and
% the count is that step's share of a full turn, rounded to a multiple of four.
%
% Rounding to four is what makes the ring usable as a port set. It puts a port
% exactly on each axis, and the axes are where Odum's flows meet a symbol --
% energy in at 180 degrees and out at 0, dissipation at 270, feedback at 90 --
% so quantising those anchors onto the ring moves them nowhere. Rounding to a
% whole number of pitches instead, as this did, left the leftover fraction of
% a turn as one wide gap at the wrap, and put no port at all on three of the
% four axes: an interaction's inflow then attached half a pitch up the notch
% edge instead of at its apex.
%
% The price is that the pitch is nominal, not exact. The step is 360/count
% rather than the pitch itself, so spacing runs up to 2/count away from 3mm --
% at most a fifth on the smallest symbols, a few percent on the largest. That
% is the better trade: a coupled pair still separates by the same distance at
% every symbol to within that, whereas the old wrap gap was a whole step wrong
% on every symbol, always in the same place.
%
% Four is the floor. A symbol with fewer ports than axes cannot take the flows
% that Odum's convention requires of it.
%
% Forty-eight is the ceiling, because that is how many ring anchors each shape
% declares below (and \energese@ringmax in energese.sty, which must agree).
% Past about 23mm of effective radius the pitch asks for more, and those a
% pathway could not name: the allocator, finding no anchor for the port it
% wanted, would walk on until it found one and attach the pathway to the far
% side of the symbol. A ring that closes coarsely is much the lesser evil, and
% 48 is a multiple of four, so the ports stay on the axes. Raising the ceiling
% means declaring more anchors in every shape.
%
% Divide before multiplying by the index. pgfmath saturates at about 16384, and
% index * 57.29578 * pitch overflows that well before the division would have
% brought it back into range.
\def\energese@ringN#1{%
    \pgfmathsetmacro{\es@reff}{#1}%
    \pgfmathsetmacro{\es@step}{57.29578 * \the\energeseanchorpitch / \es@reff}%
    \pgfmathtruncatemacro{\es@count}{max(4, min(48, 4*round(90 / \es@step)))}%
}

\def\energese@ringpt#1#2#3{%
    \energese@ringN{#2}%
    % Indices past the last port repeat earlier ones exactly, rather than
    % precessing into the gaps between them.
    %
    % Parenthesised, and not `index * 360 / count': pgfmath evaluates left to
    % right and saturates near 16384, so on a symbol carrying 46 ports or more
    % the multiplication overflows before the division can bring it back, and
    % the anchor dies with `Dimension too large'.
    \pgfmathsetmacro{\es@th}{mod(#1, \es@count) * (360 / \es@count)}%
    #3%
    \pgfmathsetmacro{\es@px}{\es@rr*cos(\es@th)}%
    \pgfmathsetmacro{\es@py}{\es@rr*sin(\es@th)}%
    \pgfpoint{\es@px pt}{\es@py pt}%
}

% The port count, published as an anchor so that code outside the shape -- the
% port allocator, the test harness -- can ask for it. PGF offers no other way
% out of a shape declaration, so the count rides in a coordinate: it is the
% offset from the node's centre, in points.
%
% Which is why `ringunit' exists beside it. An anchor's coordinate goes through
% the picture's transformation like any other, so under `scale=0.85' a count of
% 12 is read back as 10.2 -- and the allocator, quantising to a ring of ten
% points that the shape had drawn as twelve, put the heat-sink port at 240
% degrees instead of 270 and dissipation left the corner of the symbol. Reading
% the count as the *ratio* of the two offsets cancels whatever transformation
% was applied to both.
\def\energese@ringcount#1{%
    \energese@ringN{#1}%
    \pgfpoint{\es@count pt}{0pt}%
}

\def\energese@ring#1#2{%
    \anchor{ringcount}{\energese@ringcount{#1}}%
    \anchor{ringunit}{\pgfpoint{1pt}{0pt}}%
    \anchor{ring0}{\energese@ringpt{0}{#1}{#2}}%
    \anchor{ring1}{\energese@ringpt{1}{#1}{#2}}%
    \anchor{ring2}{\energese@ringpt{2}{#1}{#2}}%
    \anchor{ring3}{\energese@ringpt{3}{#1}{#2}}%
    \anchor{ring4}{\energese@ringpt{4}{#1}{#2}}%
    \anchor{ring5}{\energese@ringpt{5}{#1}{#2}}%
    \anchor{ring6}{\energese@ringpt{6}{#1}{#2}}%
    \anchor{ring7}{\energese@ringpt{7}{#1}{#2}}%
    \anchor{ring8}{\energese@ringpt{8}{#1}{#2}}%
    \anchor{ring9}{\energese@ringpt{9}{#1}{#2}}%
    \anchor{ring10}{\energese@ringpt{10}{#1}{#2}}%
    \anchor{ring11}{\energese@ringpt{11}{#1}{#2}}%
    \anchor{ring12}{\energese@ringpt{12}{#1}{#2}}%
    \anchor{ring13}{\energese@ringpt{13}{#1}{#2}}%
    \anchor{ring14}{\energese@ringpt{14}{#1}{#2}}%
    \anchor{ring15}{\energese@ringpt{15}{#1}{#2}}%
    \anchor{ring16}{\energese@ringpt{16}{#1}{#2}}%
    \anchor{ring17}{\energese@ringpt{17}{#1}{#2}}%
    \anchor{ring18}{\energese@ringpt{18}{#1}{#2}}%
    \anchor{ring19}{\energese@ringpt{19}{#1}{#2}}%
    \anchor{ring20}{\energese@ringpt{20}{#1}{#2}}%
    \anchor{ring21}{\energese@ringpt{21}{#1}{#2}}%
    \anchor{ring22}{\energese@ringpt{22}{#1}{#2}}%
    \anchor{ring23}{\energese@ringpt{23}{#1}{#2}}%
    \anchor{ring24}{\energese@ringpt{24}{#1}{#2}}%
    \anchor{ring25}{\energese@ringpt{25}{#1}{#2}}%
    \anchor{ring26}{\energese@ringpt{26}{#1}{#2}}%
    \anchor{ring27}{\energese@ringpt{27}{#1}{#2}}%
    \anchor{ring28}{\energese@ringpt{28}{#1}{#2}}%
    \anchor{ring29}{\energese@ringpt{29}{#1}{#2}}%
    \anchor{ring30}{\energese@ringpt{30}{#1}{#2}}%
    \anchor{ring31}{\energese@ringpt{31}{#1}{#2}}%
    \anchor{ring32}{\energese@ringpt{32}{#1}{#2}}%
    \anchor{ring33}{\energese@ringpt{33}{#1}{#2}}%
    \anchor{ring34}{\energese@ringpt{34}{#1}{#2}}%
    \anchor{ring35}{\energese@ringpt{35}{#1}{#2}}%
    \anchor{ring36}{\energese@ringpt{36}{#1}{#2}}%
    \anchor{ring37}{\energese@ringpt{37}{#1}{#2}}%
    \anchor{ring38}{\energese@ringpt{38}{#1}{#2}}%
    \anchor{ring39}{\energese@ringpt{39}{#1}{#2}}%
    \anchor{ring40}{\energese@ringpt{40}{#1}{#2}}%
    \anchor{ring41}{\energese@ringpt{41}{#1}{#2}}%
    \anchor{ring42}{\energese@ringpt{42}{#1}{#2}}%
    \anchor{ring43}{\energese@ringpt{43}{#1}{#2}}%
    \anchor{ring44}{\energese@ringpt{44}{#1}{#2}}%
    \anchor{ring45}{\energese@ringpt{45}{#1}{#2}}%
    \anchor{ring46}{\energese@ringpt{46}{#1}{#2}}%
    \anchor{ring47}{\energese@ringpt{47}{#1}{#2}}%
}

%% Radius at angle \es@th, per shape. -----------------------------------------

% Circle.
\def\energese@src@rad{\pgfmathsetmacro{\es@rr}{\energese@src@r}}

% Tank: a circle cut by the two straight edges of its peaked top. The arc
% begins at 11.5 degrees, so those edges run through (0.97992R, 0.19937R).
\def\energese@sto@rad{%
    \pgfmathsetmacro{\es@R}{\energese@sto@r}%
    \pgfmathsetmacro{\es@rr}{1/max(1/\es@R,
        (0.80063*abs(cos(\es@th)) + 0.97992*sin(\es@th))/(0.97992*\es@R))}%
}

% Hexagon, shoulders at 0.46H: the intersection of six half-planes, which for a
% convex outline is one over the largest of their normalised projections.
\def\energese@con@rad{%
    \pgfmathsetmacro{\es@W}{\energese@con@w}%
    \pgfmathsetmacro{\es@H}{\energese@con@h}%
    \pgfmathsetmacro{\es@rr}{1/max(abs(cos(\es@th))/\es@W,
        (0.54*\es@H*abs(cos(\es@th)) + \es@W*abs(sin(\es@th)))/(\es@W*\es@H))}%
}

% Diamond.
\def\energese@trn@rad{%
    \pgfmathsetmacro{\es@W}{\energese@trn@w}%
    \pgfmathsetmacro{\es@H}{\energese@trn@h}%
    \pgfmathsetmacro{\es@rr}{1/(abs(cos(\es@th))/\es@W + abs(sin(\es@th))/\es@H)}%
}

% Triangle: flat left edge at -W, apex at +W.
\def\energese@gan@rad{%
    \pgfmathsetmacro{\es@W}{\energese@gan@w}%
    \pgfmathsetmacro{\es@H}{\energese@gan@h}%
    \pgfmathsetmacro{\es@rr}{1/max(-cos(\es@th)/\es@W,
        (\es@H*cos(\es@th) + 2*\es@W*abs(sin(\es@th)))/(\es@W*\es@H))}%
}

% Rectangle closed by a semicircular cap of radius b centred on (a,0). Where the
% ray leaves through the cap it reaches further than the body, so the larger of
% the two is the outline.
\def\energese@prd@rad{%
    \pgfmathsetmacro{\es@a}{\energese@prd@a}%
    \pgfmathsetmacro{\es@b}{\energese@prd@b}%
    \pgfmathsetmacro{\es@disc}{\es@b*\es@b - (\es@a*sin(\es@th))^2}%
    % Three cases, and the middle one is easy to miss: a ray steep enough to
    % clear the cap entirely leaves through the flat right edge, not through a
    % cap it never meets. Collapsing that case to the cap expression put the
    % point inside the symbol.
    \pgfmathsetmacro{\es@flat}{\es@a/max(0.05,abs(cos(\es@th)))}%
    \pgfmathsetmacro{\es@side}{(cos(\es@th) > 0)
        ? ((\es@disc > 0) ? (\es@a*cos(\es@th) + sqrt(max(0,\es@disc))) : \es@flat)
        : \es@flat}%
    \pgfmathsetmacro{\es@rr}{min(\es@side, \es@b/max(0.05,abs(sin(\es@th))))}%
}

% Same construction, but the right side bows out only 0.53b, which is the arc
% through (a, +/-b) and (a + 0.53b, 0): centre (a - 0.678b), radius 1.208b.
\def\energese@lop@rad{%
    \pgfmathsetmacro{\es@a}{\energese@lop@a}%
    \pgfmathsetmacro{\es@b}{\energese@lop@b}%
    \pgfmathsetmacro{\es@cx}{\es@a - 0.678*\es@b}%
    \pgfmathsetmacro{\es@rho}{1.208*\es@b}%
    \pgfmathsetmacro{\es@disc}{\es@rho*\es@rho - (\es@cx*sin(\es@th))^2}%
    \pgfmathsetmacro{\es@flat}{\es@a/max(0.05,abs(cos(\es@th)))}%
    \pgfmathsetmacro{\es@side}{(cos(\es@th) > 0)
        ? ((\es@disc > 0) ? (\es@cx*cos(\es@th) + sqrt(max(0,\es@disc))) : \es@flat)
        : \es@flat}%
    \pgfmathsetmacro{\es@rr}{min(\es@side, \es@b/max(0.05,abs(sin(\es@th))))}%
}

% Exact, despite the notch. The outline is non-convex, so the half-plane
% construction is not obviously the right one -- but the notch replaces the
% left edge rather than cutting into a region that has one, and its supporting
% line passes through both left corners. Every ray therefore leaves through
% whichever of the three constraints binds first, which is what 1/max(...)
% computes. Checked against the sampled polygon at 0.1 degree steps and four
% aspect ratios: agreement to 0.000%.
\def\energese@int@rad{%
    \pgfmathsetmacro{\es@W}{\energese@int@w}%
    \pgfmathsetmacro{\es@H}{\energese@int@h}%
    \pgfmathsetmacro{\es@n}{0.55*\es@H}%
    \pgfmathsetmacro{\es@t}{0.85*\es@H}%
    \pgfmathsetmacro{\es@rr}{1/max(abs(sin(\es@th))/\es@H,
        (\es@H*cos(\es@th) + \es@t*abs(sin(\es@th)))/(\es@W*\es@H),
        (-\es@H*cos(\es@th) - \es@n*abs(sin(\es@th)))/(\es@H*(\es@W-\es@n)))}%
}

% The bowed side, solved rather than approximated.
%
% In the box-normalised frame X = x/W, Y = y/H all four sides are one curve:
% writing the top cubic in the symmetric parameter u = 2t-1 collapses it to
%
%     X(u) = 1.0875u - 0.0875u^3,   Y(u) = 0.715 + 0.285u^2,   u in [-1,1]
%
% -- the 0.715 at the axis and the corner at u = 1 both fall out of it -- and
% the other three sides are its reflections. A ray whose normalised direction
% has gradient m = min(|X'|,|Y'|)/max(|X'|,|Y'|) meets the side where
%
%     1.0875u - 0.0875u^3 = m (0.715 + 0.285u^2).
%
% Newton from u = m solves that to 2% in one step, 0.04% in two and to the
% resolution of a sampled path in three. Three is what runs here.
%
% What this replaces: box radius * (0.715 + 0.285 sin^4 2t), which was up to
% 19.6% outside the drawn outline between the axis and the corner -- the ring
% points there sat visibly off the symbol, which for a port means a pathway
% attached to nothing. A fixed-point iteration on the same equation converges
% far too slowly near the corner to be worth the same three lines.
\def\energese@swi@rad{%
    \pgfmathsetmacro{\es@W}{\energese@swi@w}%
    \pgfmathsetmacro{\es@H}{\energese@swi@h}%
    \pgfmathsetmacro{\es@nx}{abs(cos(\es@th))/\es@W}%
    \pgfmathsetmacro{\es@ny}{abs(sin(\es@th))/\es@H}%
    \pgfmathsetmacro{\es@m}{min(\es@nx,\es@ny)/max(\es@nx,\es@ny)}%
    \pgfmathsetmacro{\es@u}{\es@m}%
    \energese@swi@newton\energese@swi@newton\energese@swi@newton
    \pgfmathsetmacro{\es@ua}{0.715 + 0.285*\es@u*\es@u}%
    \pgfmathsetmacro{\es@ub}{1.0875*\es@u - 0.0875*\es@u*\es@u*\es@u}%
    % Which side the ray leaves through: past the diagonal of the normalised
    % box it is the top, before it the right, and the two are the same curve
    % with its coordinates swapped.
    \pgfmathsetmacro{\es@rr}{(\es@ny > \es@nx)
        ? veclen(\es@W*\es@ub, \es@H*\es@ua)
        : veclen(\es@W*\es@ua, \es@H*\es@ub)}%
}
\def\energese@swi@newton{%
    \pgfmathsetmacro{\es@u}{\es@u -
        (1.0875*\es@u - 0.0875*\es@u*\es@u*\es@u
         - \es@m*(0.715 + 0.285*\es@u*\es@u))
        / (1.0875 - 0.2625*\es@u*\es@u - 0.57*\es@m*\es@u)}%
}

% The standard `text' anchor: shifts the label so its centre lands on the
% node's origin, which is where every outline below is centred.
\def\energese@textanchor{%
    \pgf@x=-.5\wd\pgfnodeparttextbox%
    \pgf@y=-.5\ht\pgfnodeparttextbox%
    \advance\pgf@y by.5\dp\pgfnodeparttextbox%
}

%% -------------------------------------------------------------------------
%% Source -- a circle. Energy leaves to the right.
%% -------------------------------------------------------------------------
% The circle must circumscribe the text rectangle, so the radius is that
% rectangle's half-diagonal.
\pgfdeclareshape{energese source}{
    \saveddimen\energese@src@r{%
        \energese@measure%
        \pgfmathsetlength\energese@tmp{\energesemagnitude *
            veclen(\the\energese@dimx,\the\energese@dimy)}%
        \pgf@x=\energese@tmp%
    }
    \savedanchor\energese@src@ne{%
        \energese@measure%
        \pgfmathsetlength\energese@tmp{\energesemagnitude *
            veclen(\the\energese@dimx,\the\energese@dimy)}%
        \pgf@x=\energese@tmp%
        \pgf@y=\energese@tmp%
    }
    \anchor{center}{\pgfpointorigin}
    \anchor{text}{\energese@textanchor}
    \anchor{north}{\pgfpoint{0pt}{\energese@src@r}}
    \anchor{south}{\pgfpoint{0pt}{-\energese@src@r}}
    \anchor{east}{\pgfpoint{\energese@src@r}{0pt}}
    \anchor{west}{\pgfpoint{-\energese@src@r}{0pt}}
    \anchor{energy_in}{\pgfpoint{-\energese@src@r}{0pt}}
    \anchor{energy_out}{\pgfpoint{\energese@src@r}{0pt}}
    % 22.5 degrees above the horizontal, clear of the energy pathway.
    \anchor{money_in}{\pgfmathsetmacro{\es@r}{\energese@src@r}%
        \pgfpoint{0.924*\es@r pt}{0.383*\es@r pt}}
    \anchor{money_out}{\pgfmathsetmacro{\es@r}{\energese@src@r}%
        \pgfpoint{0.924*\es@r pt}{0.383*\es@r pt}}
    \anchor{feedback_in}{\pgfpoint{0pt}{\energese@src@r}}
    \anchor{feedback_out}{\pgfpoint{0pt}{\energese@src@r}}
    \anchor{heat_sink}{\pgfpoint{0pt}{-\energese@src@r}}
    \energese@ring{\energese@src@r}{\energese@src@rad}
    \anchorborder{%
        \pgf@xa=\pgf@x\pgf@ya=\pgf@y%
        \pgf@process{\energese@src@ne}%
        \pgfpointborderellipse{\pgfqpoint{\pgf@xa}{\pgf@ya}}{\pgfqpoint{\pgf@x}{\pgf@y}}%
    }
    \backgroundpath{
        \pgfpathcircle{\pgfpointorigin}{\energese@src@r}
    }
}

%% -------------------------------------------------------------------------
%% Storage -- a tank: a circle whose top is flattened into a shallow peak.
%% -------------------------------------------------------------------------
% Two constraints on the radius, plus the minimum-size floor. The circular part
% must enclose the text (its half-diagonal), and the flattened top -- which
% slopes from the peak at (0,R) down towards (0.98R, 0.2R), a gradient of
% 0.82 -- must clear the text's upper corners.
\pgfdeclareshape{energese storage}{
    \saveddimen\energese@sto@r{%
        \energese@measure%
        \pgfmathsetlength\energese@tmp{\energesemagnitude *
            max(veclen(\the\energese@dimx,\the\energese@dimy),
                \the\energese@dimy + 0.82*\the\energese@dimx)}%
        \pgf@x=\energese@tmp%
    }
    \savedanchor\energese@sto@ne{%
        \energese@measure%
        \pgfmathsetlength\energese@tmp{\energesemagnitude *
            max(veclen(\the\energese@dimx,\the\energese@dimy),
                \the\energese@dimy + 0.82*\the\energese@dimx)}%
        \pgf@x=\energese@tmp%
        \pgf@y=\energese@tmp%
    }
    \anchor{center}{\pgfpointorigin}
    \anchor{text}{\energese@textanchor}
    \anchor{north}{\pgfpoint{0pt}{\energese@sto@r}}
    \anchor{south}{\pgfpoint{0pt}{-\energese@sto@r}}
    \anchor{east}{\pgfpoint{\energese@sto@r}{0pt}}
    \anchor{west}{\pgfpoint{-\energese@sto@r}{0pt}}
    \anchor{energy_in}{\pgfpoint{-\energese@sto@r}{0pt}}
    \anchor{energy_out}{\pgfpoint{\energese@sto@r}{0pt}}
    \anchor{money_in}{\pgfmathsetmacro{\es@r}{\energese@sto@r}%
        \pgfpoint{0.94*\es@r pt}{0.34*\es@r pt}}
    \anchor{money_out}{\pgfmathsetmacro{\es@r}{\energese@sto@r}%
        \pgfpoint{-0.94*\es@r pt}{0.34*\es@r pt}}
    \anchor{feedback_in}{\pgfmathsetmacro{\es@r}{\energese@sto@r}%
        \pgfpoint{0.4*\es@r pt}{0.68*\es@r pt}}
    \anchor{feedback_out}{\pgfmathsetmacro{\es@r}{\energese@sto@r}%
        \pgfpoint{-0.4*\es@r pt}{0.68*\es@r pt}}
    \anchor{heat_sink}{\pgfpoint{0pt}{-\energese@sto@r}}
    \energese@ring{\energese@sto@r}{\energese@sto@rad}
    \anchorborder{%
        \pgf@xa=\pgf@x\pgf@ya=\pgf@y%
        \pgf@process{\energese@sto@ne}%
        \pgfpointborderellipse{\pgfqpoint{\pgf@xa}{\pgf@ya}}{\pgfqpoint{\pgf@x}{\pgf@y}}%
    }
    \backgroundpath{
        \pgfmathsetmacro{\es@r}{\energese@sto@r}
        \pgfmathsetmacro{\es@sx}{0.97992*\es@r}
        \pgfmathsetmacro{\es@sy}{0.19937*\es@r}
        \pgfpathmoveto{\pgfpoint{-\es@sx pt}{\es@sy pt}}
        \pgfpathlineto{\pgfpoint{0pt}{\es@r pt}}
        \pgfpathlineto{\pgfpoint{\es@sx pt}{\es@sy pt}}
        \pgfpatharc{11.5}{-191.5}{\es@r pt}
        \pgfpathclose
    }
}

%% -------------------------------------------------------------------------
%% Producer -- a rectangle closed by a semicircle on the right.
%% -------------------------------------------------------------------------
% The text rectangle is [-a,a] x [-b,b] and the cap is a half-circle of radius b
% centred on (a,0), so the symbol spans -a to a+b horizontally. It is
% deliberately not centred on the origin: Odum's producer is asymmetric, and the
% label belongs in the rectangular body rather than straddling the cap.
\pgfdeclareshape{energese producer}{
    % The cap occupies half the height, so the body's half-width is the
    % outline's less half the outline height.
    \saveddimen\energese@prd@a{\energese@fit{1.726}{0.710}{1.0}%
        \pgf@x=\energese@outw \advance\pgf@x by-.5\energese@outh}
    \saveddimen\energese@prd@b{\energese@fit{1.726}{0.710}{1.0}\pgf@x=\energese@outh}
    \savedanchor\energese@prd@ne{\energese@fit{1.726}{0.710}{1.0}%
        \pgf@x=\energese@outw \pgf@y=\energese@outh}
    \anchor{center}{\pgfpointorigin}
    \anchor{text}{\energese@textanchor}
    \anchor{north}{\pgfpoint{0pt}{\energese@prd@b}}
    \anchor{south}{\pgfpoint{0pt}{-\energese@prd@b}}
    \anchor{west}{\pgfpoint{-\energese@prd@a}{0pt}}
    \anchor{east}{\pgfpoint{\energese@prd@a + \energese@prd@b}{0pt}}
    \anchor{energy_in}{\pgfpoint{-\energese@prd@a}{0pt}}
    \anchor{energy_out}{\pgfpoint{\energese@prd@a + \energese@prd@b}{0pt}}
    % 45 degrees around the cap, so money leaves clear of the energy pathway.
    \anchor{money_in}{\pgfmathsetmacro{\es@a}{\energese@prd@a}%
        \pgfmathsetmacro{\es@b}{\energese@prd@b}%
        \pgfpoint{\es@a pt + 0.707*\es@b pt}{0.707*\es@b pt}}
    \anchor{money_out}{\pgfmathsetmacro{\es@b}{\energese@prd@b}%
        \pgfpoint{-\energese@prd@a}{0.5*\es@b pt}}
    \anchor{feedback_in}{\pgfmathsetmacro{\es@a}{\energese@prd@a}%
        \pgfpoint{0.3*\es@a pt}{\energese@prd@b}}
    \anchor{feedback_out}{\pgfmathsetmacro{\es@a}{\energese@prd@a}%
        \pgfpoint{-0.3*\es@a pt}{\energese@prd@b}}
    \anchor{heat_sink}{\pgfpoint{0pt}{-\energese@prd@b}}
    \energese@ring{\energese@prd@b}{\energese@prd@rad}
    \anchorborder{%
        \pgf@xa=\pgf@x\pgf@ya=\pgf@y%
        \pgf@process{\energese@prd@ne}%
        \pgfpointborderellipse{\pgfqpoint{\pgf@xa}{\pgf@ya}}{\pgfqpoint{\pgf@x}{\pgf@y}}%
    }
    \backgroundpath{
        \pgfmathsetmacro{\es@a}{\energese@prd@a}
        \pgfmathsetmacro{\es@b}{\energese@prd@b}
        \pgfpathmoveto{\pgfpoint{-\es@a pt}{\es@b pt}}
        \pgfpathlineto{\pgfpoint{\es@a pt}{\es@b pt}}
        \pgfpatharc{90}{-90}{\es@b pt}
        \pgfpathlineto{\pgfpoint{-\es@a pt}{-\es@b pt}}
        \pgfpathclose
    }
}

%% -------------------------------------------------------------------------
%% Consumer -- a hexagon, pointed top and bottom, with vertical sides.
%% -------------------------------------------------------------------------
% Measured from reference-images/complete-odum-symbols.png by flood-filling the
% symbol interior: aspect 0.877, with the shoulders at 0.46 of the half-height,
% and energy entering and leaving through the vertical sides.
%
% Let H = lambda*b. Requiring the label's corner at y = b to sit inside the
% sloped edge gives W >= a*0.54*lambda/(lambda-1); requiring the unlabelled
% symbol to hit Odum's 0.877 aspect gives W = 0.877*lambda*b. Solving the two
% together at a = b yields lambda = 1.616, at which both coefficients coincide
% at 1.417 -- so H = 1.616b and W = 1.417*max(a,b) fits every label exactly and
% keeps the canonical proportion.
%
% Deriving lambda rather than picking it matters: an earlier conservative rule
% (H = b/0.46) also fitted every label, but made the consumer 2.17 times the
% height of a producer carrying the same text. Odum's own diagrams put that
% ratio at 1.69, which is what 1.616 delivers.
\pgfdeclareshape{energese consumer}{
    \saveddimen\energese@con@w{\energese@fit{0.877}{1.0}{0.46}\pgf@x=\energese@outw}
    % A wide label drives W up, which would flatten the hexagon well past
    % anything Odum draws (his labelled consumers measure 1.09-1.12 against a
    % bare 0.877). Grow H alongside W to hold the aspect at 1.08 or below --
    % a single-line label at 1.15 leaves the hexagon looking stretched even
    % though a two-line label at the same ratio does not.
    \saveddimen\energese@con@h{\energese@fit{0.877}{1.0}{0.46}\pgf@x=\energese@outh}
    \savedanchor\energese@con@ne{\energese@fit{0.877}{1.0}{0.46}%
        \pgf@x=\energese@outw \pgf@y=\energese@outh}
    \anchor{center}{\pgfpointorigin}
    \anchor{text}{\energese@textanchor}
    \anchor{north}{\pgfpoint{0pt}{\energese@con@h}}
    \anchor{south}{\pgfpoint{0pt}{-\energese@con@h}}
    \anchor{east}{\pgfpoint{\energese@con@w}{0pt}}
    \anchor{west}{\pgfpoint{-\energese@con@w}{0pt}}
    \anchor{energy_in}{\pgfpoint{-\energese@con@w}{0pt}}
    \anchor{energy_out}{\pgfpoint{\energese@con@w}{0pt}}
    % Money runs through the vertical sides, above the energy pathway.
    \anchor{money_in}{\pgfmathsetmacro{\es@w}{\energese@con@w}%
        \pgfmathsetmacro{\es@h}{\energese@con@h}%
        \pgfpoint{\es@w pt}{0.3*\es@h pt}}
    \anchor{money_out}{\pgfmathsetmacro{\es@w}{\energese@con@w}%
        \pgfmathsetmacro{\es@h}{\energese@con@h}%
        \pgfpoint{-\es@w pt}{0.3*\es@h pt}}
    % Feedback arrives on the upper slopes, as in Odum's originals.
    \anchor{feedback_in}{\pgfmathsetmacro{\es@w}{\energese@con@w}%
        \pgfmathsetmacro{\es@h}{\energese@con@h}%
        \pgfpoint{0.5*\es@w pt}{0.73*\es@h pt}}
    \anchor{feedback_out}{\pgfmathsetmacro{\es@w}{\energese@con@w}%
        \pgfmathsetmacro{\es@h}{\energese@con@h}%
        \pgfpoint{-0.5*\es@w pt}{0.73*\es@h pt}}
    \anchor{heat_sink}{\pgfpoint{0pt}{-\energese@con@h}}
    \energese@ring{\energese@con@h}{\energese@con@rad}
    \anchorborder{%
        \pgf@xa=\pgf@x\pgf@ya=\pgf@y%
        \pgf@process{\energese@con@ne}%
        \pgfpointborderellipse{\pgfqpoint{\pgf@xa}{\pgf@ya}}{\pgfqpoint{\pgf@x}{\pgf@y}}%
    }
    \backgroundpath{
        \pgfmathsetmacro{\es@w}{\energese@con@w}
        \pgfmathsetmacro{\es@h}{\energese@con@h}
        \pgfmathsetmacro{\es@sh}{0.46*\es@h}
        \pgfpathmoveto{\pgfpoint{0pt}{\es@h pt}}
        \pgfpathlineto{\pgfpoint{\es@w pt}{\es@sh pt}}
        \pgfpathlineto{\pgfpoint{\es@w pt}{-\es@sh pt}}
        \pgfpathlineto{\pgfpoint{0pt}{-\es@h pt}}
        \pgfpathlineto{\pgfpoint{-\es@w pt}{-\es@sh pt}}
        \pgfpathlineto{\pgfpoint{-\es@w pt}{\es@sh pt}}
        \pgfpathclose
    }
}

%% -------------------------------------------------------------------------
%% Transaction -- Odum's exchange diamond.
%% -------------------------------------------------------------------------
% W = a + 2.3b and H = b + a/2 keep the text's corners just inside the edges
% while giving an unlabelled diamond the aspect of 2.23 measured from Odum's
% sheet. Energy runs through the lower half, money through the upper.
\pgfdeclareshape{energese transaction}{
    \saveddimen\energese@trn@w{\energese@fit{2.226}{0.5}{0.5}\pgf@x=\energese@outw}
    \saveddimen\energese@trn@h{\energese@fit{2.226}{0.5}{0.5}\pgf@x=\energese@outh}
    \savedanchor\energese@trn@ne{\energese@fit{2.226}{0.5}{0.5}%
        \pgf@x=\energese@outw \pgf@y=\energese@outh}
    \anchor{center}{\pgfpointorigin}
    \anchor{text}{\energese@textanchor}
    \anchor{north}{\pgfpoint{0pt}{\energese@trn@h}}
    \anchor{south}{\pgfpoint{0pt}{-\energese@trn@h}}
    \anchor{east}{\pgfpoint{\energese@trn@w}{0pt}}
    \anchor{west}{\pgfpoint{-\energese@trn@w}{0pt}}
    \anchor{energy_in}{\pgfmathsetmacro{\es@w}{\energese@trn@w}%
        \pgfmathsetmacro{\es@h}{\energese@trn@h}%
        \pgfpoint{-0.5*\es@w pt}{-0.5*\es@h pt}}
    \anchor{energy_out}{\pgfmathsetmacro{\es@w}{\energese@trn@w}%
        \pgfmathsetmacro{\es@h}{\energese@trn@h}%
        \pgfpoint{0.5*\es@w pt}{-0.5*\es@h pt}}
    \anchor{money_in}{\pgfmathsetmacro{\es@w}{\energese@trn@w}%
        \pgfmathsetmacro{\es@h}{\energese@trn@h}%
        \pgfpoint{0.5*\es@w pt}{0.5*\es@h pt}}
    \anchor{money_out}{\pgfmathsetmacro{\es@w}{\energese@trn@w}%
        \pgfmathsetmacro{\es@h}{\energese@trn@h}%
        \pgfpoint{-0.5*\es@w pt}{0.5*\es@h pt}}
    \anchor{feedback_in}{\pgfmathsetmacro{\es@w}{\energese@trn@w}%
        \pgfmathsetmacro{\es@h}{\energese@trn@h}%
        \pgfpoint{0.5*\es@w pt}{0.5*\es@h pt}}
    \anchor{feedback_out}{\pgfmathsetmacro{\es@w}{\energese@trn@w}%
        \pgfmathsetmacro{\es@h}{\energese@trn@h}%
        \pgfpoint{-0.5*\es@w pt}{0.5*\es@h pt}}
    \anchor{heat_sink}{\pgfpoint{0pt}{-\energese@trn@h}}
    \energese@ring{\energese@trn@h}{\energese@trn@rad}
    \anchorborder{%
        \pgf@xa=\pgf@x\pgf@ya=\pgf@y%
        \pgf@process{\energese@trn@ne}%
        \pgfpointborderellipse{\pgfqpoint{\pgf@xa}{\pgf@ya}}{\pgfqpoint{\pgf@x}{\pgf@y}}%
    }
    \backgroundpath{
        \pgfmathsetmacro{\es@w}{\energese@trn@w}
        \pgfmathsetmacro{\es@h}{\energese@trn@h}
        \pgfpathmoveto{\pgfpoint{0pt}{\es@h pt}}
        \pgfpathlineto{\pgfpoint{\es@w pt}{0pt}}
        \pgfpathlineto{\pgfpoint{0pt}{-\es@h pt}}
        \pgfpathlineto{\pgfpoint{-\es@w pt}{0pt}}
        \pgfpathclose
    }
}

%% -------------------------------------------------------------------------
%% Interaction -- pointed on the right, notched on the left.
%% -------------------------------------------------------------------------
% The notch is where two or more inflows converge, so `energy_in' sits at its
% apex rather than out on the bounding box. Half-height H = 1.15b, and
% W = a + 1.17b gives an unlabelled symbol the aspect of 1.89 measured from
% Odum's sheet. The notch depth and right taper both scale with H, and W clears
% whichever of the two intrudes further.
%
% Taper and notch are deep (0.85H and 0.55H) so the front point and back notch
% come to sharp angles as Odum draws them. They shorten the flat top and bottom
% without touching the bounding box, so the aspect above is unaffected.
\pgfdeclareshape{energese interaction}{
    \saveddimen\energese@int@w{\energese@fit{1.886}{0.708}{0.85}\pgf@x=\energese@outw}
    \saveddimen\energese@int@h{\energese@fit{1.886}{0.708}{0.85}\pgf@x=\energese@outh}
    \savedanchor\energese@int@ne{\energese@fit{1.886}{0.708}{0.85}%
        \pgf@x=\energese@outw \pgf@y=\energese@outh}
    \anchor{center}{\pgfpointorigin}
    \anchor{text}{\energese@textanchor}
    \anchor{north}{\pgfpoint{0pt}{\energese@int@h}}
    \anchor{south}{\pgfpoint{0pt}{-\energese@int@h}}
    \anchor{east}{\pgfpoint{\energese@int@w}{0pt}}
    % Due west of an interaction is inside the notch, where there is no
    % outline: -W on the axis is the void between the two upper and lower left
    % corners, and a pathway drawn to it starts in mid-air. The westmost point
    % of the symbol itself is the notch apex.
    \anchor{west}{\pgfmathsetmacro{\es@w}{\energese@int@w}%
        \pgfmathsetmacro{\es@h}{\energese@int@h}%
        \pgfpoint{-\es@w pt + 0.55*\es@h pt}{0pt}}
    % The notch apex, where converging inflows meet.
    \anchor{energy_in}{\pgfmathsetmacro{\es@w}{\energese@int@w}%
        \pgfmathsetmacro{\es@h}{\energese@int@h}%
        \pgfpoint{-\es@w pt + 0.55*\es@h pt}{0pt}}
    \anchor{energy_out}{\pgfpoint{\energese@int@w}{0pt}}
    \anchor{money_in}{\pgfmathsetmacro{\es@w}{\energese@int@w}%
        \pgfmathsetmacro{\es@h}{\energese@int@h}%
        \pgfpoint{\es@w pt - 0.35*\es@h pt}{0.5*\es@h pt}}
    \anchor{money_out}{\pgfmathsetmacro{\es@w}{\energese@int@w}%
        \pgfmathsetmacro{\es@h}{\energese@int@h}%
        \pgfpoint{-\es@w pt + 0.275*\es@h pt}{0.5*\es@h pt}}
    \anchor{feedback_in}{\pgfpoint{0pt}{\energese@int@h}}
    \anchor{feedback_out}{\pgfmathsetmacro{\es@w}{\energese@int@w}%
        \pgfpoint{-0.3*\es@w pt}{\energese@int@h}}
    \anchor{heat_sink}{\pgfpoint{0pt}{-\energese@int@h}}
    \energese@ring{\energese@int@h}{\energese@int@rad}
    \anchorborder{%
        \pgf@xa=\pgf@x\pgf@ya=\pgf@y%
        \pgf@process{\energese@int@ne}%
        \pgfpointborderellipse{\pgfqpoint{\pgf@xa}{\pgf@ya}}{\pgfqpoint{\pgf@x}{\pgf@y}}%
    }
    \backgroundpath{
        \pgfmathsetmacro{\es@w}{\energese@int@w}
        \pgfmathsetmacro{\es@h}{\energese@int@h}
        \pgfmathsetmacro{\es@notch}{0.55*\es@h}
        \pgfmathsetmacro{\es@taper}{0.85*\es@h}
        \pgfpathmoveto{\pgfpoint{-\es@w pt}{\es@h pt}}
        \pgfpathlineto{\pgfpoint{\es@w pt - \es@taper pt}{\es@h pt}}
        \pgfpathlineto{\pgfpoint{\es@w pt}{0pt}}
        \pgfpathlineto{\pgfpoint{\es@w pt - \es@taper pt}{-\es@h pt}}
        \pgfpathlineto{\pgfpoint{-\es@w pt}{-\es@h pt}}
        \pgfpathlineto{\pgfpoint{-\es@w pt + \es@notch pt}{0pt}}
        \pgfpathclose
    }
}

%% -------------------------------------------------------------------------
%% Gain -- Odum's constant-gain amplifier: a triangle pointing right.
%% -------------------------------------------------------------------------
% Measured aspect 0.923 against a reference of 0.917. H = 2.6b and
% W = 2.4*max(a,b) hold that proportion.
%
% A triangle is a poor container for centred text: fitting the full corners of a
% label's bounding box would force H = 3.09b, making the symbol three times a
% producer's height. Instead the label's middle band is guaranteed to fit (at
% y = b/2 the edge is at 1.48a) and the extreme corners of the bounding box --
% normally whitespace above the x-height -- may sit outside. Odum labels
% amplifiers with a single character; long labels here will look wrong.
\pgfdeclareshape{energese gain}{
    \saveddimen\energese@gan@w{\energese@fit{0.917}{0.45}{0.55}\pgf@x=\energese@outw}
    \saveddimen\energese@gan@h{\energese@fit{0.917}{0.45}{0.55}\pgf@x=\energese@outh}
    \savedanchor\energese@gan@ne{\energese@fit{0.917}{0.45}{0.55}%
        \pgf@x=\energese@outw \pgf@y=\energese@outh}
    \anchor{center}{\pgfpointorigin}
    \anchor{text}{\energese@textanchor}
    \anchor{north}{\pgfmathsetmacro{\es@h}{\energese@gan@h}\pgfpoint{0pt}{0.5*\es@h pt}}
    \anchor{south}{\pgfmathsetmacro{\es@h}{\energese@gan@h}\pgfpoint{0pt}{-0.5*\es@h pt}}
    \anchor{east}{\pgfpoint{\energese@gan@w}{0pt}}
    \anchor{west}{\pgfpoint{-\energese@gan@w}{0pt}}
    \anchor{energy_in}{\pgfpoint{-\energese@gan@w}{0pt}}
    \anchor{energy_out}{\pgfpoint{\energese@gan@w}{0pt}}
    \anchor{money_in}{\pgfmathsetmacro{\es@w}{\energese@gan@w}%
        \pgfmathsetmacro{\es@h}{\energese@gan@h}%
        \pgfpoint{0.3*\es@w pt}{0.35*\es@h pt}}
    \anchor{money_out}{\pgfmathsetmacro{\es@h}{\energese@gan@h}%
        \pgfpoint{-\energese@gan@w}{0.5*\es@h pt}}
    \anchor{feedback_in}{\pgfmathsetmacro{\es@h}{\energese@gan@h}\pgfpoint{0pt}{0.5*\es@h pt}}
    \anchor{heat_sink}{\pgfmathsetmacro{\es@h}{\energese@gan@h}\pgfpoint{0pt}{-0.5*\es@h pt}}
    \energese@ring{\energese@gan@h}{\energese@gan@rad}
    \anchorborder{%
        \pgf@xa=\pgf@x\pgf@ya=\pgf@y%
        \pgf@process{\energese@gan@ne}%
        \pgfpointborderellipse{\pgfqpoint{\pgf@xa}{\pgf@ya}}{\pgfqpoint{\pgf@x}{\pgf@y}}%
    }
    \backgroundpath{
        \pgfmathsetmacro{\es@w}{\energese@gan@w}
        \pgfmathsetmacro{\es@h}{\energese@gan@h}
        \pgfpathmoveto{\pgfpoint{-\es@w pt}{\es@h pt}}
        \pgfpathlineto{\pgfpoint{\es@w pt}{0pt}}
        \pgfpathlineto{\pgfpoint{-\es@w pt}{-\es@h pt}}
        \pgfpathclose
    }
}

%% -------------------------------------------------------------------------
%% Loop-limited converter -- a rectangle whose right side bows outward.
%% -------------------------------------------------------------------------
% Reference interior 54 x 53, aspect 1.019, of which the flat body is about
% 40 x 53 and the bulge adds 14 -- so the bulge is 0.53b and the body is taller
% than wide. The style therefore sets a minimum width below the minimum height;
% that ratio, not the path, is what produces the reference aspect when
% unlabelled. The bulge is a cubic whose control offset of 1.333d puts the
% midpoint at exactly a + d.
\pgfdeclareshape{energese loop_limited}{
    \saveddimen\energese@lop@a{\energese@fit{1.019}{0.740}{1.0}%
        \pgf@x=\energese@outw \advance\pgf@x by-.265\energese@outh}
    \saveddimen\energese@lop@b{\energese@fit{1.019}{0.740}{1.0}\pgf@x=\energese@outh}
    \savedanchor\energese@lop@ne{\energese@fit{1.019}{0.740}{1.0}%
        \pgf@x=\energese@outw \pgf@y=\energese@outh}
    \anchor{center}{\pgfpointorigin}
    \anchor{text}{\energese@textanchor}
    \anchor{north}{\pgfpoint{0pt}{\energese@lop@b}}
    \anchor{south}{\pgfpoint{0pt}{-\energese@lop@b}}
    \anchor{west}{\pgfpoint{-\energese@lop@a}{0pt}}
    \anchor{east}{\pgfmathsetmacro{\es@a}{\energese@lop@a}%
        \pgfmathsetmacro{\es@b}{\energese@lop@b}%
        \pgfpoint{\es@a pt + 0.53*\es@b pt}{0pt}}
    \anchor{energy_in}{\pgfpoint{-\energese@lop@a}{0pt}}
    \anchor{energy_out}{\pgfmathsetmacro{\es@a}{\energese@lop@a}%
        \pgfmathsetmacro{\es@b}{\energese@lop@b}%
        \pgfpoint{\es@a pt + 0.53*\es@b pt}{0pt}}
    \anchor{money_in}{\pgfmathsetmacro{\es@a}{\energese@lop@a}%
        \pgfmathsetmacro{\es@b}{\energese@lop@b}%
        \pgfpoint{\es@a pt + 0.36*\es@b pt}{0.5*\es@b pt}}
    \anchor{money_out}{\pgfmathsetmacro{\es@b}{\energese@lop@b}%
        \pgfpoint{-\energese@lop@a}{0.5*\es@b pt}}
    \anchor{feedback_in}{\pgfmathsetmacro{\es@a}{\energese@lop@a}%
        \pgfpoint{0.3*\es@a pt}{\energese@lop@b}}
    \anchor{feedback_out}{\pgfmathsetmacro{\es@a}{\energese@lop@a}%
        \pgfpoint{-0.3*\es@a pt}{\energese@lop@b}}
    \anchor{heat_sink}{\pgfpoint{0pt}{-\energese@lop@b}}
    \energese@ring{\energese@lop@b}{\energese@lop@rad}
    \anchorborder{%
        \pgf@xa=\pgf@x\pgf@ya=\pgf@y%
        \pgf@process{\energese@lop@ne}%
        \pgfpointborderellipse{\pgfqpoint{\pgf@xa}{\pgf@ya}}{\pgfqpoint{\pgf@x}{\pgf@y}}%
    }
    \backgroundpath{
        \pgfmathsetmacro{\es@a}{\energese@lop@a}
        \pgfmathsetmacro{\es@b}{\energese@lop@b}
        \pgfmathsetmacro{\es@k}{0.707*\es@b}
        \pgfpathmoveto{\pgfpoint{-\es@a pt}{\es@b pt}}
        \pgfpathlineto{\pgfpoint{\es@a pt}{\es@b pt}}
        \pgfpathcurveto{\pgfpoint{\es@a pt + \es@k pt}{\es@b pt}}
                       {\pgfpoint{\es@a pt + \es@k pt}{-\es@b pt}}
                       {\pgfpoint{\es@a pt}{-\es@b pt}}
        \pgfpathlineto{\pgfpoint{-\es@a pt}{-\es@b pt}}
        \pgfpathclose
    }
}

%% -------------------------------------------------------------------------
%% Switch -- on/off switching action: a square with all four sides bowed in.
%% -------------------------------------------------------------------------
% Reference aspect 1.036, met by W = 1.55*max(a,b) and H = 1.496b. Control
% points at 0.62 of the half-extent put each bowed side's midpoint at 0.715 of
% the half-extent, and that is where a label's edge has least room: 1.399 is the
% coefficient at which the label exactly touches, so 1.55 leaves 12% clearance.
\pgfdeclareshape{energese switch}{
    \saveddimen\energese@swi@w{\energese@fit{1.036}{0.715}{0.715}\pgf@x=\energese@outw}
    \saveddimen\energese@swi@h{\energese@fit{1.036}{0.715}{0.715}\pgf@x=\energese@outh}
    \savedanchor\energese@swi@ne{\energese@fit{1.036}{0.715}{0.715}%
        \pgf@x=\energese@outw \pgf@y=\energese@outh}
    \anchor{center}{\pgfpointorigin}
    \anchor{text}{\energese@textanchor}
    \anchor{north}{\pgfmathsetmacro{\es@h}{\energese@swi@h}\pgfpoint{0pt}{0.715*\es@h pt}}
    \anchor{south}{\pgfmathsetmacro{\es@h}{\energese@swi@h}\pgfpoint{0pt}{-0.715*\es@h pt}}
    \anchor{east}{\pgfmathsetmacro{\es@w}{\energese@swi@w}\pgfpoint{0.715*\es@w pt}{0pt}}
    \anchor{west}{\pgfmathsetmacro{\es@w}{\energese@swi@w}\pgfpoint{-0.715*\es@w pt}{0pt}}
    \anchor{energy_in}{\pgfmathsetmacro{\es@w}{\energese@swi@w}\pgfpoint{-0.715*\es@w pt}{0pt}}
    \anchor{energy_out}{\pgfmathsetmacro{\es@w}{\energese@swi@w}\pgfpoint{0.715*\es@w pt}{0pt}}
    \anchor{money_in}{\pgfmathsetmacro{\es@w}{\energese@swi@w}%
        \pgfmathsetmacro{\es@h}{\energese@swi@h}%
        \pgfpoint{0.78*\es@w pt}{0.45*\es@h pt}}
    \anchor{money_out}{\pgfmathsetmacro{\es@w}{\energese@swi@w}%
        \pgfmathsetmacro{\es@h}{\energese@swi@h}%
        \pgfpoint{-0.78*\es@w pt}{0.45*\es@h pt}}
    \anchor{feedback_in}{\pgfmathsetmacro{\es@h}{\energese@swi@h}\pgfpoint{0pt}{0.715*\es@h pt}}
    \anchor{feedback_out}{\pgfmathsetmacro{\es@w}{\energese@swi@w}%
        \pgfmathsetmacro{\es@h}{\energese@swi@h}%
        \pgfpoint{-0.4*\es@w pt}{0.79*\es@h pt}}
    \anchor{heat_sink}{\pgfmathsetmacro{\es@h}{\energese@swi@h}\pgfpoint{0pt}{-0.715*\es@h pt}}
    \energese@ring{\energese@swi@h}{\energese@swi@rad}
    \anchorborder{%
        \pgf@xa=\pgf@x\pgf@ya=\pgf@y%
        \pgf@process{\energese@swi@ne}%
        \pgfpointborderellipse{\pgfqpoint{\pgf@xa}{\pgf@ya}}{\pgfqpoint{\pgf@x}{\pgf@y}}%
    }
    \backgroundpath{
        \pgfmathsetmacro{\es@w}{\energese@swi@w}
        \pgfmathsetmacro{\es@h}{\energese@swi@h}
        \pgfmathsetmacro{\es@cw}{0.62*\es@w}
        \pgfmathsetmacro{\es@ch}{0.62*\es@h}
        \pgfmathsetmacro{\es@qw}{0.45*\es@w}
        \pgfmathsetmacro{\es@qh}{0.45*\es@h}
        \pgfpathmoveto{\pgfpoint{-\es@w pt}{\es@h pt}}
        \pgfpathcurveto{\pgfpoint{-\es@qw pt}{\es@ch pt}}{\pgfpoint{\es@qw pt}{\es@ch pt}}
                       {\pgfpoint{\es@w pt}{\es@h pt}}
        \pgfpathcurveto{\pgfpoint{\es@cw pt}{\es@qh pt}}{\pgfpoint{\es@cw pt}{-\es@qh pt}}
                       {\pgfpoint{\es@w pt}{-\es@h pt}}
        \pgfpathcurveto{\pgfpoint{\es@qw pt}{-\es@ch pt}}{\pgfpoint{-\es@qw pt}{-\es@ch pt}}
                       {\pgfpoint{-\es@w pt}{-\es@h pt}}
        \pgfpathcurveto{\pgfpoint{-\es@cw pt}{-\es@qh pt}}{\pgfpoint{-\es@cw pt}{\es@qh pt}}
                       {\pgfpoint{-\es@w pt}{\es@h pt}}
        \pgfpathclose
    }
}

%% -------------------------------------------------------------------------
%% Ground -- the heat-sink terminator. A fixed-size glyph, never labelled.
%% -------------------------------------------------------------------------
\pgfdeclareshape{energese ground}{
    \savedanchor\centerpoint{\pgfpointorigin}
    \anchor{center}{\centerpoint}
    \anchor{north}{\pgfpoint{0pt}{0pt}}
    \anchor{south}{\pgfpoint{0pt}{-0.3cm}}
    \anchor{east}{\pgfpoint{0.2cm}{-0.1cm}}
    \anchor{west}{\pgfpoint{-0.2cm}{-0.1cm}}
    \backgroundpath{
        \pgfpathmoveto{\pgfpoint{0pt}{0pt}}
        \pgfpathlineto{\pgfpoint{0pt}{-0.1cm}}
        \pgfpathmoveto{\pgfpoint{-0.2cm}{-0.1cm}}
        \pgfpathlineto{\pgfpoint{0.2cm}{-0.1cm}}
        \pgfpathmoveto{\pgfpoint{-0.1cm}{-0.17cm}}
        \pgfpathlineto{\pgfpoint{0.1cm}{-0.17cm}}
        \pgfpathmoveto{\pgfpoint{-0.05cm}{-0.24cm}}
        \pgfpathlineto{\pgfpoint{0.05cm}{-0.24cm}}
        \pgfusepath{draw}
    }
}
